\documentclass[a4paper]{article}

% Basic packages
% Español (México) con XeLaTeX: fontspec para fuentes Unicode, polyglossia
% para la división silábica y los nombres del documento en la variante mexicana.
\usepackage{fontspec}
\usepackage{polyglossia}
\setdefaultlanguage[variant=mexican]{spanish}
% Los nombres chinos van como «pinyin (original)»: xeCJK da a los caracteres
% Han su propia fuente; sin ella XeLaTeX los omite con un aviso.
% FandolSong no cubre algunos caracteres de nombres (U+73FA); WenQuanYi Zen Hei sí.
\usepackage[AutoFallBack=true]{xeCJK}
\setCJKmainfont{FandolSong-Regular.otf}
\setCJKfallbackfamilyfont{\CJKrmdefault}{WenQuanYi Zen Hei}
% polyglossia no traduce los nombres de listings.
\AtBeginDocument{\providecommand{\lstlistingname}{}\renewcommand{\lstlistingname}{Listado}%
  \providecommand{\lstlistlistingname}{}\renewcommand{\lstlistlistingname}{Índice de listados}}
\usepackage{amsmath, amssymb}
\usepackage{graphicx}
\usepackage[margin=2.5cm]{geometry}
\usepackage[most]{tcolorbox}
\usepackage{etoolbox}
\usepackage{listings}
\usepackage{hyperref}
\usepackage{booktabs}
\usepackage{subcaption}
\usepackage{float}

% Highlight boxes
\newtcolorbox{knowledgebox}[1]{enhanced, colback=blue!5!white, colframe=blue!75!black, colbacktitle=blue!75!black, coltitle=white, fonttitle=\bfseries, title={#1}, attach boxed title to top left={yshift=-2mm, xshift=2mm}, boxrule=1pt, sharp corners}
\newtcolorbox{importantbox}[1]{enhanced, colback=yellow!10!white, colframe=yellow!80!black, colbacktitle=yellow!80!black, coltitle=black, fonttitle=\bfseries, title={#1}, sharp corners}
\newtcolorbox{warningbox}[1]{enhanced, colback=red!5!white, colframe=red!75!black, colbacktitle=red!75!black, coltitle=white, fonttitle=\bfseries, title={#1}, sharp corners}

% Listings
\lstset{language=Python, basicstyle=\ttfamily\small, keywordstyle=\color{blue}, stringstyle=\color{red!60!black}, commentstyle=\color{green!60!black}, breaklines=true, frame=single, numbers=left, numberstyle=\tiny\color{gray}, captionpos=b, extendedchars=true}

% Front-page metadata
\newcommand{\notetitle}{[Berkeley LLM Agents F24] Agents for Software Development — Graham Neubig}
\newcommand{\noteauthors}{Elaboradas a partir de materiales públicos del curso}
\newcommand{\notedate}{14 de octubre de 2024}
\newcommand{\videochannel}{Berkeley RDI}
\newcommand{\videopublishdate}{14 de octubre de 2024}
\newcommand{\videoduration}{1:01:21}
\newcommand{\videourl}{https://www.youtube.com/live/f9L9Fkq-8K4}
\newcommand{\videocoverpath}{cover.jpg}
\newcommand{\repourl}{https://github.com/hqhq1025/ai-course-notes}

% Footer with repo info
\usepackage{fancyhdr}
\pagestyle{fancy}
\fancyhf{}
\fancyfoot[C]{\thepage}
\fancyfoot[R]{\footnotesize\color{gray}\href{\repourl}{AI Course Notes}}
\renewcommand{\headrulewidth}{0pt}
\renewcommand{\footrulewidth}{0pt}

\begin{document}

\begin{titlepage}
\centering
{\Large Notas de la entrevista\par}
\vspace{1.2cm}
{\huge\bfseries \notetitle\par}
\vspace{0.8cm}
{\large \noteauthors\par}
\vspace{0.3cm}
{\large \notedate\par}
\vspace{1.2cm}

\ifdefempty{\videocoverpath}{
{\small Al generar las notas, indique la ruta local de la portada del video.\par}
}{
\includegraphics[width=0.82\textwidth,height=0.45\textheight,keepaspectratio]{\videocoverpath}\par
}

\vfill
\begin{tcolorbox}[width=0.9\textwidth, colback=black!2!white, colframe=black!60, sharp corners]
\textbf{Autor o canal del video}: \videochannel\par
\textbf{Fecha de publicación}: \videopublishdate\par
\textbf{Duración del video}: \videoduration\par
\textbf{Enlace del video}: \href{\videourl}{\nolinkurl{\videourl}}\par
\textbf{Página del proyecto}: \href{\repourl}{\nolinkurl{\repourl}}\par
\end{tcolorbox}
\end{titlepage}

\tableofcontents
\newpage

\section{Introducción: por qué están en auge los agentes de desarrollo de software}
\subsection{Contexto de la industria y plataforma}
Graham Neubig abre presentándose como profesor de la CMU y científico en jefe de All Hands AI, y relata su acumulación doble entre el desarrollo de software y la inteligencia artificial. All Hands AI mantiene el framework de código abierto OpenHands (antes OpenDev), cuyo objetivo es empaquetar la capacidad de «escribir código» en un flujo lineal que un agente pueda ejecutar, de modo que los desarrolladores puedan concentrarse en los requerimientos y la revisión, y dejarle al agente las pruebas, la búsqueda y los parches repetitivos.

Señala que una plataforma abierta permite que la comunidad construya en conjunto herramientas de depuración y repositorios de ejemplo, extraiga entradas de issues reales, dé seguimiento a las operaciones del agente en repositorios públicos y retroalimente las mejoras al modelo. Este ciclo cerrado permite que la investigación llegue directamente a la práctica de ingeniería, en lugar de quedarse en los puntos de referencia de laboratorio.

\subsection{Empoderar al público general}
Menciona en particular el artículo de Mark Andreessen de 2011, ``Why Software Is Eating the World'', que señala que todas las industrias están siendo redefinidas por el software: del entretenimiento a la agricultura, la defensa y los servicios médicos, el software ha ido sustituyendo procesos físicos, y el ámbito del que se encargará el software en el futuro solo va a crecer. Esto también explica por qué es necesario que más personas cuenten con la capacidad de construir software rápidamente, en lugar de que el desarrollo de software siga monopolizado por un pequeño número de ingenieros profesionales.

Al presentar la motivación, Graham usa como analogía varios avances científicos recientes impulsados por sistemas de software, para subrayar que el progreso científico contemporáneo está profundamente entrelazado con el código. Para él, el valor del agente está en poder transferir la capacidad de ingeniería a un grupo más amplio de investigadores, gerentes de producto y equipos interdisciplinarios, de modo que «crear software» se vuelva una capacidad generalizada en lugar de un recurso escaso.

\begin{knowledgebox}{La predicción de Mark Andreessen}
El artículo de 2011 ``Why Software Is Eating the World'' sostiene que el software se convertirá en la nueva infraestructura sobre la que se apoyan todos los servicios de la industria, y que la capacidad de automatizar la construcción de software será el factor decisivo para la competitividad. Toda la charla de Graham busca llevar esta idea a la práctica en la ingeniería de agentes.
\end{knowledgebox}

\begin{importantbox}{La propuesta de valor del agente}
La promesa central de los agentes de desarrollo de software es convertir un sistema multimodal con capacidad de comprender, modificar, verificar y retroalimentar código en un ejecutor capaz de iterar un producto de manera continua, igual que un desarrollador humano.
\end{importantbox}

\subsection{Resumen de la sección}
\begin{itemize}
    \item El software ya domina la mayoría de las industrias, y la escasez de capacidad de construcción impulsa la necesidad real de agentes.
    \item Graham, al contrastar el Premio Nobel con el ámbito del software, nos recuerda que crear software ya es una de las competencias centrales de la ciencia humana.
\end{itemize}
\section{La plataforma OpenHands y la arquitectura del Agent}
\subsection{Componentes de la arquitectura de OpenHands}
OpenHands ofrece una plataforma de desarrollo de Agent contenedorizada y orquestable: el Agent puede operar simultáneamente la terminal, el editor y el navegador dentro de un entorno sandbox controlado; el backend puede conectarse a varios LLM (como code-davinci y modelos de código abierto), y a través de una cadena de herramientas de webhook e intérpretes de pruebas completa el ciclo cerrado de entrada→inferencia→acción.

Los componentes se comunican entre sí mediante un bus de eventos y una cola de mensajes, lo que permite sustituir el modelo o las herramientas en tiempo de ejecución. La plataforma también define metadatos para cada acción (forma en que se desencadena, contexto, salida esperada), lo que facilita la reversión y la depuración.

\subsection{Interfaz multimodal y auditoría}
\begin{figure}[H]
\centering
\includegraphics[width=\textwidth]{frame-01-openhands.jpg}
\caption{La pantalla de OpenHands muestra el flujo en cascada del Agent: percepción del entorno, inferencia del LLM y ejecución de herramientas. \protect\footnotemark}
\end{figure}
\footnotetext{Intervalo de tiempo en el video: 00:04:45--00:04:50.}

La plataforma pone énfasis en la «observabilidad»: cada acción de cada Agent (comandos de terminal, clics en el navegador, modificaciones de archivos) queda registrada, y los resultados de las pruebas y los registros fluyen de vuelta a un panel de auditoría central. OpenHands también ofrece interfaces para que una persona pueda interrumpir el proceso, usar la vista de depuración o insertar scripts auxiliares; esta capacidad de tipo plug-and-play hace que el Agent deje de ser una caja negra.

\begin{knowledgebox}{Capacidades centrales de OpenHands}
OpenHands incluye un adaptador multimodal, un motor de decisión basado en el LLM, ejecutores de herramientas (shell de Docker, pruebas, navegador) y un panel de auditoría visual. Cada parte se comunica mediante un bus de eventos, y el modelo o las herramientas pueden sustituirse en tiempo de ejecución.
\end{knowledgebox}

\begin{warningbox}{Observabilidad y límites de seguridad}
El Agent registra cada comando, pero también debe restringir las operaciones con privilegios elevados y garantizar que el sandbox pueda restaurarse; de lo contrario, los registros de desarrollo y el estado del entorno pueden filtrar información sensible y volverse difíciles de reproducir.
\end{warningbox}

\subsection{Resumen de la sección}
\begin{itemize}
    \item OpenHands conecta la percepción, la inferencia, la acción y la auditoría del Agent en un pipeline modular.
    \item Los registros transparentes y las interfaces de intervención humana son la garantía clave que hace que los equipos de ingeniería estén dispuestos a poner en producción un Agent.
\end{itemize}
\section{SWE-bench: un banco de pruebas para la evaluación de la ingeniería de software}
\subsection{La evolución de Lite y Full}
Graham presenta a continuación SWE-bench: un marco de evaluación en el que un Agente resuelve issues reales de GitHub. Cada muestra incluye el repositorio original, la descripción del issue y la suite de pruebas; la tarea del Agente es comprender la semántica y el contexto, escribir un parche y pasar las pruebas de ese repositorio para verificar su corrección.

SWE-bench se divide en dos categorías: Lite (300 problemas relativamente simples) y Full (2,294 issues de repositorios reales). Lite se concentra en comprender comandos, parches y pruebas básicas; Full extiende el problema a dependencias entre varios archivos, llamadas entre módulos y entornos complejos. El Agente está limitado a una ventana de contexto, por lo que necesita diseñar estrategias eficaces de prompt, recuperación por índice y razonamiento local.

\subsection{Proceso de evaluación y métricas}
\begin{figure}[H]
\centering
\includegraphics[width=\textwidth]{frame-02-swebench.jpg}
\caption{Panel de estadísticas de evaluación de SWE-bench, que muestra la tasa de éxito de los conjuntos de datos Lite y Full. \protect\footnotemark}
\end{figure}
\footnotetext{Intervalo de tiempo en el video: 00:15:20--00:15:25.}

El proceso de evaluación consta de tres pasos: 1) el Agente lee el issue y recupera el contexto relevante en el repositorio; 2) genera y aplica un parche (edita archivos, ejecuta \texttt{git diff}); 3) ejecuta las pruebas objetivo y archiva los resultados junto con los registros de salida. El sistema calcula la tasa de éxito según \texttt{pass/fail} y retroalimenta las muestras fallidas al proceso de entrenamiento para mejorar el prompt y las estrategias de herramientas.

\begin{table}[H]
\centering
\begin{tabular}{p{0.35\textwidth} p{0.58\textwidth}}
\toprule
Dimensión & Descripción \\
\midrule
Tipo de tarea & Extraer el objetivo del issue, ubicar los archivos relevantes, generar un parche y ejecutar las pruebas correspondientes. \\
\midrule
Lite & 300 puntos propensos a error, orientados a verificar la estructura del prompt y la capacidad de corrección rápida; los mejores Agentes actuales alcanzan una tasa de éxito de entre 40\%--50\%. \\
\midrule
Full & 2,294 desafíos, incluidas llamadas asíncronas, migración de dependencias y parches de seguridad; la tasa de éxito general actual es inferior a 20\%. \\
\bottomrule
\end{tabular}
\caption{Dimensiones de diseño de SWE-bench y resultados actuales.}
\end{table}

\begin{knowledgebox}{Énfasis de la evaluación de SWE-bench}
SWE-bench enfatiza el ciclo completo de ingeniería: el Agente debe entregar el patch, ejecutar las pruebas y proporcionar los registros; no basta con dar el diff. La división entre Lite y Full permite a los equipos de investigación optimizar primero las capacidades básicas antes de extenderse a escenarios complejos.
\end{knowledgebox}

\begin{importantbox}{Indicadores cuantitativos}
Graham enfatiza: una tasa de éxito de 40\% en el nivel Lite ya significa que el Agente puede sustituir parte de las tareas repetitivas, pero la baja tasa de éxito en el nivel Full indica que todavía estamos en una era de «semiautomatización».
\end{importantbox}

\begin{warningbox}{Límites de cobertura}
SWE-bench evalúa principalmente la codificación y las pruebas; no incluye tareas más amplias de ingeniería de software como el diseño, la revisión de código, el despliegue y el mantenimiento, de modo que las capacidades evaluadas siguen inclinadas hacia problemas de tipo corrección.
\end{warningbox}

\subsection{Resumen de la sección}
\begin{itemize}
    \item SWE-bench ofrece un espectro doble que va de lo simple a lo complejo, y le da al Agente un objetivo que puede alcanzarse de manera gradual.
    \item Los indicadores actuales siguen muy por debajo de los de un ingeniero humano, y todavía se necesitan nuevas evaluaciones para etapas más amplias del ciclo de vida del software.
\end{itemize}
\section{Práctica de ingeniería con agentes: planeación, ejecución y verificación}
\subsection{Detalles de Plan-Act-Verify}
En la ingeniería real, el agente debe contar con un ciclo de acción explícito: primero leer el issue, invocar el servicio de recuperación para fijar el contexto, planear el parche, ejecutar la modificación del código, ejecutar las pruebas y, después, ajustar con la retroalimentación de los registros. Graham enfatiza la combinación de varias herramientas: el LLM se enfoca en generar, los scripts se encargan de operar, el marco de pruebas verifica los resultados y, por último, la evidencia se sube al dashboard para la revisión humana.

Señala que en cada ronda de Plan primero debe generarse una lista de tareas, enumerando de manera explícita qué archivos se modificarán y qué pruebas deben ejecutarse. La etapa de Act, en cambio, se vincula a scripts y comandos de terminal concretos, de modo que el agente escriba el código y lo ejecute realmente; la etapa de Verify ejecuta las pruebas y el análisis estático, y las muestras que fallan disparan una nueva Plan.

\subsection{Colaboración humano-máquina y cobertura de herramientas}
\begin{figure}[H]
\centering
\includegraphics[width=\textwidth]{frame-03-human-loop.jpg}
\caption{El supervisor humano puede pausar, revisar las decisiones e intervenir mientras el agente está en ejecución, corrigiendo así errores en tiempo real. \protect\footnotemark}
\end{figure}
\footnotetext{Intervalo de tiempo en el video: 00:32:05--00:32:10.}

OpenHands permite al desarrollador insertar una revisión humana en cualquier paso: cuando el agente planea ejecutar un comando peligroso (como un despliegue o una escritura en la base de datos), puede pausar el proceso, revisar los registros y, de ser necesario, modificar el prompt o ejecutar manualmente parte de los comandos. Las tareas de larga duración también envían checkpoints, lo que facilita el seguimiento del estado.

\begin{importantbox}{Mejor práctica}
Se implementa un pipeline para el agente: primero plan (segmentar la tarea, encontrar las dependencias), luego act (edición de código, compilación, pruebas) y, por último, verify (registros, instantáneas, capacidad de reversión). Los pasos que fallan deben rastrearse hacia atrás y pasarle al LLM el contexto clave.
\end{importantbox}

\begin{knowledgebox}{El valor de la intervención humana}
Permitir que el ser humano observe la conducta del agente durante su ejecución, y pausarlo antes de que realice una operación peligrosa, es clave para mejorar la confiabilidad, sobre todo en compilaciones de publicación u operaciones de base de datos.
\end{knowledgebox}

\begin{warningbox}{Ventana de contexto y alucinaciones}
El contexto del agente es limitado: si se introduce de una vez toda la estructura del repositorio, se rebasa la ventana legible, por lo que primero debe hacerse la extracción o recuperación de información; de lo contrario, el LLM producirá parches ilusorios por falta de información del entorno.
\end{warningbox}

\subsection{Resumen de la sección}
\begin{itemize}
    \item El ciclo planeación-ejecución-verificación hace que el agente sea controlable dentro del flujo de ingeniería.
    \item La colaboración humano-máquina (pausar, revisar registros, saltar de contexto) es un medio necesario para evitar errores graves.
\end{itemize}
\section{Direcciones futuras y desafíos de gobernanza}
\subsection{Entrenamiento de agentes y datos}
Graham señala que el entrenamiento actual de agentes todavía no aprovecha suficientemente los datos de «estilo agente» (agent chains, autorretroalimentación, corrección de fallas); en el futuro se necesitarán más agent-structured corpora para que el modelo aprenda a planificar y a corregir. Otra dirección es reforzar la evaluación human-in-the-loop: mediante una revisión conjunta entre humanos y agentes, el modelo aprende a detectar y corregir sus propios errores.

Agrega que el Agent training debería incorporar los registros del proceso como parte de los datos de entrenamiento, para que el modelo no solo vea el patch final, sino también el Plan, los pasos de ejecución y los resultados de Verify.

\subsection{Gobernanza y seguridad}
\begin{knowledgebox}{Núcleo de la investigación futura}
El énfasis está en: entrenar el comportamiento integrado del agente, incorporar trayectorias de corrección humana y construir evaluaciones más granulares (por ejemplo, si los registros del proceso están alineados y si el patch es mantenible).
\end{knowledgebox}

\begin{warningbox}{Gobernanza y seguridad}
La automatización excesiva de los agentes puede provocar decisiones de alto riesgo: por ejemplo, desplegar un patch directamente sin supervisión humana o invocar APIs sensibles entre repositorios; para llevarlo a producción es necesario introducir mecanismos de aprobación, auditoría y reversión.
\end{warningbox}

\begin{importantbox}{Desglose de direcciones}
Las tres rutas que recomienda Graham: 1) entrenar con agent-specific data, 2) construir o ampliar evaluaciones del tipo SWE-bench, 3) incorporar retroalimentación human-in-the-loop en tiempo real y mecanismos de seguridad.
\end{importantbox}

\subsection{Resumen de la sección}
\begin{itemize}
    \item Resolver el vacío en el entrenamiento y la evaluación de agentes es el punto de partida para que los agentes lleguen a una etapa de madurez.
    \item La seguridad, la aprobación y la intervención humana no deben considerarse añadidos, sino parte integral del ciclo de vida del agente.
\end{itemize}
\section{Caso de ingeniería: repositorios reales con participación de Agent}
\subsection{Commits semiautomatizados}
Graham compartió su práctica con los repositorios: en algunos proyectos, «la mitad del contenido de los commits lo redacta primero un Agent de codificación, y después él lo revisa y corrige». El Agent primero lee el Issue, encuentra los archivos relevantes mediante resumen y búsqueda, después redacta un parche preliminar y, por último, un ingeniero humano lo integra.

Él prevé que, en el próximo año o dos, los Agent se encargarán de más tareas cotidianas y rutinarias, para que los ingenieros puedan concentrar su atención en el diseño de alto nivel. Por ahora, el Agent puede manejar sin problemas los scripts de Package con dependencias compartidas, las operaciones básicas de CLI y las correcciones simples de errores (bug fix); los escenarios complejos todavía requieren revisión humana.

\subsection{Revisión del proceso y retroalimentación}
Después de cada ejecución del Agent, el equipo guarda un registro final: incluye los archivos obtenidos mediante fetch de forma proactiva, los comandos invocados, los resultados de las pruebas y las causas de los fracasos. Estos registros se convierten en muestras de entrenamiento y recordatorios de retrollamada, para que el Agent aprenda a incorporar en la cadena de prompts los pasos útiles de las muestras fallidas.

También recalcó: «Observar el proceso de ejecución del Agent y dar correcciones inmediatas mejora más la calidad del modelo que puntuar únicamente el resultado.» Esta retroalimentación basada en la observación incluye pausas humanas, la inyección de prompts y hacer que el Agent vuelva a intentarlo después de corregirlo.

\begin{knowledgebox}{Experiencia del caso}
Graham ya usa el Agent para redactar en varios repositorios, y los commits que aporta el Agent representan cerca de la mitad del volumen de trabajo. Esto demuestra que el Agent ya tiene un valor práctico en las tareas repetitivas.
\end{knowledgebox}

\begin{importantbox}{Recomendación práctica}
Despliegue primero al Agent en tareas de menos de 10 minutos (como el formateo, los parches o las actualizaciones de dependencias), arme un proceso de review humano y, después, extiéndalo poco a poco hacia longer-running job.
\end{importantbox}

\begin{warningbox}{No deje que el Agent se despliegue solo}
Aun cuando el Agent pueda completar ciertos cambios, no se le debe dejar llevarlos solo a producción: deben conservarse la revisión humana, las pruebas de regresión y la aprobación (approve) manual.
\end{warningbox}

\subsection{Resumen de la sección}
\begin{itemize}
    \item Los experimentos con el Agent en repositorios reales ya muestran el potencial de la semiautomatización, pero todavía no bastan para sustituir por completo a los ingenieros humanos.
    \item Llevar un registro del proceso y corregir de inmediato son las prácticas fundamentales para mejorar la confiabilidad del Agent.
\end{itemize}
\section{Observabilidad y prácticas de depuración}
\subsection{Registro y pipeline de monitoreo}
Para que cada ejecución del Agent se pueda rastrear, OpenHands abstrae la terminal, el navegador, las pruebas y los registros en un flujo de eventos unificado. Cada evento incluye una marca de tiempo, la herramienta invocada, una instantánea del contexto y el resultado esperado, lo que facilita a los ingenieros localizar el punto de la falla.

\begin{itemize}
    \item Registro de terminal: registra cada comando de shell y su código de retorno, y se sube junto con los resultados de las pruebas.
    \item Interacción de navegador y red: captura clics, envíos de formularios y solicitudes/respuestas, lo que facilita reproducir sesiones HTTP y flujos de API.
    \item Diff de archivos/código: guarda automáticamente los fragmentos de diff modificados por el Agent, las rutas de los archivos relacionados y el contexto.
\end{itemize}

\begin{table}[H]
\centering
\begin{tabular}{p{0.25\textwidth} p{0.35\textwidth} p{0.33\textwidth}}
\toprule
Nivel & Uso & Ejemplo \\
\midrule
Nivel de comando & Captura cada operación de shell y su resultado & \texttt{pip install}, \texttt{pytest tests/foo.py} \\
\midrule
Nivel de contexto & Registra la versión del contexto usada y los datos recuperados & commit de Git + Issue summary correspondientes \\
\midrule
Nivel de resultado & Archiva si las pruebas pasaron o fallaron, análisis de diff & Unit test failed with stack trace \\
\bottomrule
\end{tabular}
\caption{Niveles de observación y usos del registro en OpenHands.}
\end{table}

\subsection{Modo de depuración y reversión rápida}
Cuando el Agent planea ejecutar un comando de alto riesgo (como una publicación o una migración de base de datos), el sistema entra automáticamente en modo de depuración: pausa el proceso, captura la instantánea actual y le pide a un humano que confirme el comando. Si la ejecución del comando falla, la plataforma puede revertirla automáticamente con base en los registros y el diff, evitando contaminar la rama principal.

\begin{knowledgebox}{Recomendación de depuración}
Hacer que el Agent registre en el log la «intención» y el «resultado esperado» de cada paso puede acelerar la revisión manual y formar muestras de entrenamiento.
\end{knowledgebox}

\begin{warningbox}{El uso indebido de los registros puede filtrar información sensible}
Los registros de depuración pueden contener claves de API, rutas y trazas de pila, por lo que es necesario establecer anonimización y control de acceso para evitar filtraciones de privacidad durante las auditorías.
\end{warningbox}

\subsection{Resumen de la sección}
\begin{itemize}
    \item Con un flujo de eventos unificado y registros con plantilla se puede acotar rápidamente el alcance de las fallas.
    \item El modo de depuración ofrece una doble garantía: confirmación manual y reversión automática.
\end{itemize}
\section{Métricas del agente y controles de calidad}
\subsection{Métricas clave}
Además de la tasa de éxito en SWE-bench, es necesario prestar atención a otras métricas: la tasa de intervención humana, el tiempo de espera promedio, la cobertura de pruebas y la tasa de diferencias. Estas métricas ayudan al equipo a determinar si el agente realmente reduce la carga de ingeniería en lugar de agregar costos adicionales.

\begin{importantbox}{Desglose de métricas}
Graham recomienda prestar atención a cuatro dimensiones: 1) tasa de éxito (aprobado/reprobado), 2) tasa de intervención humana (cuántas veces se requiere pausar), 3) tiempo de ejecución (si está dentro de un rango aceptable), 4) interpretabilidad (si la salida ofrece una explicación razonable).
\end{importantbox}

\subsection{Nodos de control de calidad}
En el pipeline del agente, los controles de calidad pueden colocarse en los nodos Plan, Act o Verify: Plan requiere una revisión humana de los pasos generados, la etapa Act necesita una revisión de los comandos por ejecutar y la etapa Verify consiste en pruebas y análisis estático. Solo cuando se superan todos los controles se permite enviar el patch.

\begin{table}[H]
\centering
\begin{tabular}{p{0.3\textwidth} p{0.33\textwidth} p{0.33\textwidth}}
\toprule
Nodo de control & Criterio & Respuesta \\
\midrule
Plan & Si los pasos del plan son válidos y si implican operaciones con permisos & Revisión humana o adición de un prompt \\
\midrule
Act & Si el comando involucra recursos sensibles (bases de datos, servicios en la nube) & Pausar y solicitar aprobación \\
\midrule
Verify & Si las pruebas o el análisis estático pasan en su totalidad & Revertir, agregar pruebas o someter a revisión humana \\
\bottomrule
\end{tabular}
\caption{Disposición típica de los controles de calidad en el flujo del agente.}
\end{table}

\begin{warningbox}{No te fijes solo en la tasa de éxito}
Una tasa de éxito alta puede significar que el agente está limitado a tareas de bajo riesgo; en un entorno de producción real también es necesario considerar de manera conjunta la evaluación humana y el costo de ejecución.
\end{warningbox}

\subsection{Resumen de la sección}
\begin{itemize}
    \item La tasa de éxito es solo una de las métricas; hay que evaluarla junto con la tasa de intervención, el tiempo de ejecución y otros factores.
    \item Los controles de calidad ayudan a intervenir en los nodos clave, evitando que el agente publique directamente e introduzca riesgos.
\end{itemize}
\section{Glosario y lista de herramientas}
\subsection{Términos comunes}
Para facilitar la colaboración del equipo, Graham recomienda mantener un glosario: Plan-Act-Verify, Agent Chain, Observability Stream, entre otros; un lenguaje unificado reduce los malentendidos.

\begin{itemize}
    \item Plan-Act-Verify: divide con claridad los pasos de planeación, ejecución y verificación del agente.
    \item Agent Chain: combina varios agentes en un flujo de ejecución encadenado para tareas complejas.
    \item Observability Stream: el flujo de datos unificado de los registros, comandos y resultados de pruebas que produce el agente.
\end{itemize}

\subsection{Lista de herramientas}
OpenHands incorpora varios conjuntos de herramientas: Terminal Runner, Browser Controller, Test Executor, Diff Reporter. Al desplegar, se puede mantener un entorno y permisos independientes para cada una, y registrar su uso.

\begin{table}[H]
\centering
\begin{tabular}{p{0.3\textwidth} p{0.6\textwidth}}
\toprule
Herramienta & Uso \\
\midrule
Terminal Runner & Ejecuta comandos de shell y registra stdout/stderr. \\
\midrule
Browser Controller & Realiza interacción con páginas, extrae datos y los reescribe en el entorno. \\
\midrule
Test Executor & Ejecuta pytest, npm test, etcétera, y recopila los resultados. \\
\midrule
Diff Reporter & Reúne los patches y genera una descripción y un diff en markdown. \\
\bottomrule
\end{tabular}
\caption{Principales herramientas de OpenHands y su función.}
\end{table}

\begin{knowledgebox}{Estrategia de combinación de herramientas}
Procura usar herramientas con permisos distintos en pasos diferentes para reducir el riesgo de operaciones erróneas; por ejemplo, limita las operaciones de base de datos a un comando independiente y establece una lógica de reversión.
\end{knowledgebox}

\subsection{Resumen de la sección}
\begin{itemize}
    \item Un glosario facilita la colaboración entre equipos y evita malentendidos sobre la intención del agente.
    \item Una lista de herramientas ayuda a modularizar distintos tipos de operaciones, facilitando la asignación de permisos y la auditoría.
\end{itemize}
\section{Colaboración de ingeniería y flujo de cambios}
\subsection{Perspectiva de gestión de proyectos}
Los equipos que trabajan junto con un Agent necesitan redefinir el «estado de la tarea»: los patches que genera el Agent deben incorporarse a una lista de revisión pendiente y mantenerse alineados con el backlog, los issues y la descomposición de tareas. Graham recomienda usar el enfoque de «los humanos se encargan de los requisitos + el Agent se encarga de la ejecución» para dividir los requisitos en varias tareas pequeñas, y luego registrar la actividad del Agent en el sprint board.

\begin{itemize}
    \item Tratar la fase de planeación del Agent como una «tarea de desarrollo» revisada por el PM o el tech lead.
    \item En la fase Act del Agent, conservar un mecanismo de confirmación humana que garantice que solo se ejecuten los comandos aprobados con los permisos correspondientes.
    \item En la fase Verify, publicar el reporte de pruebas en el issue para que QA lo revise en la regresión.
\end{itemize}

\subsection{Entrega y comunicación de riesgos}
Cuando el Agent ayuda a implementar un patch, el estado de la entrega debe identificarse claramente como «borrador del Agent + confirmación humana». Si el Agent falla, hay que notificar automáticamente al owner correspondiente, con el registro y el diff adjuntos, para que pueda hacerse cargo de inmediato.

 \begin{knowledgebox}{Matriz de comunicación}
Registro rápido: el Agent falla → se notifica al owner; el Agent tiene éxito pero requiere revisión → se da seguimiento al PR; el Agent ejecuta un comando peligroso → se solicita aprobación.
\end{knowledgebox}

\begin{warningbox}{No tratar al Agent como una caja negra}
En la comunicación hay que explicar qué está haciendo el Agent y por qué lo hace, para reducir el temor del equipo a confiar en él y garantizar que, en caso de falla, alguien pueda tomar el relevo con rapidez.
\end{warningbox}

\subsection{Resumen de la sección}
\begin{itemize}
    \item La gestión de proyectos necesita tratar la salida del Agent como un elemento de trabajo de primera clase e incorporarlo al tablero.
    \item La comunicación de riesgos y el mecanismo de aprobación son requisitos previos para el despliegue del Agent a gran escala.
\end{itemize}

\section{Resumen rápido del contenido de la entrevista}
\subsection{0:00--15:00: Contexto y plataforma}
La parte inicial destaca el contexto de All Hands AI y OpenHands, con énfasis en delegar la capacidad del agente a un equipo más amplio. Graham relató su filosofía de mantener su GitHub «verde»: quiere ver cada vez más commits generados en un primer borrador por el agente y después refinados por humanos.

\begin{itemize}
    \item 0--15 minutos: contexto y posicionamiento de la plataforma.
    \item 15--30 minutos: diseño de SWE-bench y datos experimentales.
    \item 30--45 minutos: casos prácticos, intervención humana y depuración.
    \item 45--60 minutos: hoja de ruta de investigación futura y gobernanza.
\end{itemize}

\subsection{Resumen de la sección}
\begin{itemize}
    \item La entrevista se desarrolla siguiendo una línea de tiempo, y abarca motivación, punto de referencia, práctica y futuro.
    \item Cada parte subraya el hilo conductor de que «el agente necesita la participación humana».
\end{itemize}

\section{Prácticas de despliegue y monitoreo}
\subsection{Integración con CI/CD}
Graham sugiere combinar el Agent con SWE-bench dentro del pipeline de CI: ejecutar el conjunto Lite todos los días o en cada commit, y generar un issue con las muestras que fallan para que el equipo de modelos las estudie. Los scripts automatizados pueden ejecutar Lite/Full en paralelo en varios entornos, y subir los registros para que las personas los revisen.

\begin{lstlisting}[language=bash]
#!/usr/bin/env bash
set -ex
cd /workspace/swe-bench
python run_swe_bench.py --split lite --model codegen
python run_swe_bench.py --split full --tests smoke
tar czf swe-bench-results.tar.gz outputs/
gzip -c outputs/logs/latest.log > swe-bench-results/logs/latest.log.gz
notify-team --subject "SWE-bench run finished" --files swe-bench-results.tar.gz
\end{lstlisting}

El \texttt{notify-team} del script puede conectarse a Slack o correo electrónico, y enviar junto el diff de la falla y las entradas/salidas generales.

\begin{knowledgebox}{Estrategia de ejecución de CI/CD}
Divide la ejecución del Agent en «precalentamiento, ejecución y reproducción»: en la etapa de precalentamiento se reconstruyen el repositorio y las dependencias, en la etapa de ejecución se ejecuta el Agent, y en la etapa de reproducción se suben los registros y el diff para que los desarrolladores los revisen.
\end{knowledgebox}

\subsection{Experimentos a pequeña escala}
Antes de ampliar el despliegue del Agent, primero realiza 1 o 2 tareas piloto en el repositorio \texttt{sandbox}: usa el Agent para generar el patch, pero revísalo solo en la red interna, y mantén cada ronda de ejecución en menos de 15 minutos para asegurar que sea controlable.

\begin{importantbox}{Recomendaciones para experimentos de corto plazo}
Conserva puntos de verificación humanos: permite que el Agent modifique menos de 3 archivos, y toda operación debe generar un patch en \texttt{git} que una persona confirme después.
\end{importantbox}

\begin{warningbox}{Evitar la contaminación de la rama principal}
No dejes que el Agent haga push directamente a `main`/`master`; debe generar el patch en una rama feature y fusionarla hasta que termine la revisión.
\end{warningbox}

\subsection{Resumen de la sección}
\begin{itemize}
    \item Integrar el Agent en el CI permite detectar regressions de manera periódica, y retroalimentar los trace de las fallas.
    \item Los experimentos a pequeña escala en la red interna pueden reducir el riesgo de iteración y acumular experiencia.
\end{itemize}

\section{Recomendaciones de la entrevista y preguntas frecuentes}
\subsection{Repaso de las preguntas del público}
Graham respondió varias preguntas frecuentes: ¿puede el Agent encargarse de tareas que no son de código? ¿cómo se mide la tasa de éxito del Agent? ¿cuándo hace falta que un humano tome el control? Sus respuestas reflejan la convicción central de la «colaboración humano-máquina»: el Agent se encarga de la parte repetible y orquestable, y el humano conserva la autoridad de revisión y control de riesgos.

\begin{itemize}
    \item «¿Puede el Agent hacer análisis de requisitos?» → Puede ayudar a organizar el texto, pero el PM todavía tiene que indagar los puntos ciegos.
    \item «¿Cómo se toma el control cuando algo falla?» → El diff que genera el Agent lo revisa un human in the loop, y si no se aprueba se revierte.
    \item «¿Cómo se conserva la traceability?» → OpenHands registra todas las acciones, y el registro incluye la marca de tiempo, el comando y el resultado.
\end{itemize}

\subsection{Lista de recomendaciones}
\begin{table}[H]
\centering
\begin{tabular}{p{0.33\textwidth} p{0.63\textwidth}}
\toprule
Tema & Recomendación \\
\midrule
Alineación de requisitos & Dividir los requisitos en tareas pequeñas y especificar la salida esperada durante la etapa de Plan. \\
\midrule
Habilitación de registros & Hacer que el Agent escriba «intención + resultado» en el registro, para facilitar el análisis humano. \\
\midrule
Control de seguridad & Configurar las operaciones de base de datos y los comandos de despliegue como pasos que requieren approve humano. \\
\bottomrule
\end{tabular}
\caption{Panorama rápido de las recomendaciones frecuentes.}
\end{table}

\subsection{Resumen de la sección}
\begin{itemize}
    \item El Q\&A de la entrevista vuelve a subrayar que el Agent no sustituye al ingeniero humano, sino que es su compañero de apoyo.
    \item Una lista de recomendaciones clara facilita que el equipo lleve a la práctica el contenido de la entrevista.
\end{itemize}

\section{Percepción de datos y panel de evaluación}
\subsection{Vistas centrales del panel}
Graham recomienda construir un conjunto de paneles de evaluación de Agentes que monitoreen en tiempo real la tasa de éxito, el número de intervenciones, la duración media de ejecución y los registros de activación de las compuertas de seguridad. Cada indicador del panel debería vincularse con el log del evento correspondiente, para poder rastrear la causa raíz cuando un indicador se comporte de manera anómala.

\begin{table}[H]
\centering
\begin{tabular}{p{0.28\textwidth} p{0.42\textwidth} p{0.28\textwidth}}
\toprule
Indicador & Descripción & Objetivo \\
\midrule
Tasa de aprobación & Proporción en que el Agente ejecuta un patch y pasa las pruebas & \(\geq 40\%\) (Lite), \(\geq 20\%\) (Full) \\
\midrule
Tasa de intervención humana & Proporción que requiere pausa o aprobación de una persona & \(\leq 10\%\) \\
\midrule
Duración de ejecución & Tiempo medio desde Plan hasta la finalización de Verify & \(\leq 20 \text{ minutos}\) \\
\midrule
Activación de seguridad & Número de comandos que involucran la base de datos o el despliegue & Requiere revisión humana \\
\bottomrule
\end{tabular}
\caption{Indicadores centrales que se vigilan en el panel del Agente.}
\end{table}

\begin{knowledgebox}{Dependencias de datos del panel}
Cada indicador debería vincularse con el log del evento (comandos, diff, salida de las pruebas), para facilitar el rastreo y entrenar modelos interpretables.
\end{knowledgebox}

\subsection{Ciclo de retroalimentación del panel}
El panel no solo debe monitorear, también debe sustentar un proceso de «retroalimentación -> ajuste»: cuando un indicador retrocede, el sistema envía automáticamente el log de la falla al equipo de Agentes, para que produzcan una nueva combinación de prompt y herramientas, y la nueva estrategia se despliegue al entorno de pruebas.

\begin{importantbox}{Práctica de retroalimentación de datos}
Diseña reportes automatizados: cuando un indicador sea anómalo (por ejemplo, tasa de aprobación < 30\%), abre automáticamente un issue, adjunta el diff y el log, y notifica al equipo, para evitar que el problema quede olvidado por mucho tiempo.
\end{importantbox}

\subsection{Resumen de la sección}
\begin{itemize}
    \item El panel en tiempo real es una ventana clave para juzgar la madurez de un Agente.
    \item Convertir los eventos anómalos en un ciclo de retroalimentación automatizado permite acelerar la iteración de la estrategia.
\end{itemize}

\section{Cultura de equipo y capacitación}
\subsection{Estrategia de implementación cultural}
Graham enfatiza: cuando el Agent no tiene un desempeño sólido, las personas dudan de su valor; por eso hay que difundir dentro del equipo la cultura de «Agent + humano», dando como ejemplo que el Agent se encarga de tareas repetitivas, mientras el humano conserva la revisión final. En cada retrospectiva de sprint hay que mostrar los resultados del Agent.

\begin{enumerate}
    \item Compartir semanalmente casos de éxito del Agent (por ejemplo, el Agent genera un diff y pasa 3 pruebas).
    \item Elaborar un documento de «Agent Etiquette» que precise cuándo se puede invocar al Agent y en qué casos es obligatoria la revisión humana.
    \item Establecer un programa de Mentor, en el que un ingeniero con experiencia guíe a los nuevos integrantes antes de que se familiaricen con las herramientas del Agent.
\end{enumerate}

\subsection{Capacitación y guías}
El material de capacitación debe incluir terminología básica (Plan-Act-Verify/Observability Stream), los scripts más utilizados y advertencias de riesgo. Al terminar la capacitación se puede organizar una práctica real: hacer que los alumnos, en un sandbox, dejen que el Agent procese un Issue, mientras observan a un lado y registran los puntos de intervención.

\begin{warningbox}{Retos comunes de la capacitación}
Los nuevos integrantes tienden a depender demasiado del Agent; se recomienda establecer, en las prácticas reales, un umbral de «revisión humana» para asegurarse de que sepan verificar la salida del Agent.
\end{warningbox}

\begin{knowledgebox}{Transmisión de conocimiento}
Registrar en documentos y videos el proceso de cada colaboración entre humanos y máquinas, para facilitar que los equipos compartan la experiencia entre sí y reducir el riesgo de que «solo una persona sepa usar el Agent».
\end{knowledgebox}

\subsection{Resumen de la sección}
\begin{itemize}
    \item La cultura, la capacitación y las prácticas reales sostienen juntas la implementación a largo plazo del Agent.
    \item Una etiquette clara y un mecanismo de retrospectiva pueden eliminar la desconfianza hacia el Agent.
\end{itemize}

\section{Gobernanza de seguridad y proceso de revisión}
\subsection{Diseño de la cadena de auditoría}
Cada comando, diff de archivos y resultado de pruebas que genera OpenHands se envía a la cadena de auditoría, y el servicio de auditoría los escribe en el middleware de eventos. Esto permite que el equipo revise fácilmente, después de los hechos, «quién hizo que el agente ejecutara qué comando», y lo contraste con la política de seguridad.

\begin{table}[H]
\centering
\begin{tabular}{p{0.25\textwidth} p{0.4\textwidth} p{0.33\textwidth}}
\toprule
Fase de auditoría & Contenido & Salida \\
\midrule
Captura de eventos & Comandos, navegación, operaciones de archivos & ID de log + marca de tiempo \\
\midrule
Alineación con la política & Verifica si se activan recursos sensibles & Warning / Approval request \\
\midrule
Archivo del informe & Genera correo/issue & Registro de auditoría reproducible \\
\bottomrule
\end{tabular}
\caption{Las tres fases del proceso de auditoría de seguridad.}
\end{table}

\begin{importantbox}{Elementos de la auditoría}
El registro debe incluir el contexto (issue, diff, pruebas), para que una persona pueda juzgar el riesgo con rapidez.
\end{importantbox}

\subsection{Aprobación y permisos}
Cuando el agente planea ejecutar un comando sensible (por ejemplo, una migración de base de datos o un despliegue en la nube), el módulo de aprobación se activa automáticamente. La decisión de aprobación incluye el comando que la desencadena, el entorno de invocación, el resultado esperado, el diff relacionado y el estado de la prueba anterior; quien aprueba solo necesita confirmar o rechazar.

\begin{warningbox}{El ritmo de la aprobación}
No se debe aprobar en exceso al punto de bloquear la iteración; pero en los escenarios de alto riesgo es necesario establecer multi-party approval, para evitar un punto único de falla.
\end{warningbox}

\subsection{Resumen de la sección}
\begin{itemize}
    \item La cadena de auditoría de OpenHands registra el comportamiento del agente, para usarse en revisiones de seguridad posteriores.
    \item Establecer la aprobación y los permisos de manera adecuada puede reforzar la seguridad sin reducir la eficiencia.
\end{itemize}
\section{Repaso práctico y plantillas de casos}
\subsection{Ejemplo de corrección de código}
Graham describió un caso típico: después de que el Agent encuentra un bug, en la etapa de Plan se escribe la lista de archivos a modificar, en la etapa de Act se usa \texttt{python apply\_patch} para generar el diff, y en la etapa de Verify se ejecuta \texttt{pytest}. Todo el proceso queda registrado en el log, y si la prueba falla se revierte automáticamente.

\begin{knowledgebox}{Ejemplo de flujo}
Plan: extraer el Issue, listar archivos y funciones; Act: ejecutar el Agent para escribir el patch y guardar el diff; Verify: ejecutar las pruebas, y si fallan, enviar el log al issue.
\end{knowledgebox}

\subsection{Registro de fallas y correcciones}
El equipo registra cada muestra de falla en una «matriz de fallas»: incluye la causa de la falla (prueba, permisos, sintaxis), la salida del Agent y el contenido de la intervención humana. Esta matriz es material importante para entrenar los prompts.

\begin{warningbox}{No borren el log de fallas}
El log de fallas contiene información semántica clave, que se usa para generar datos de entrenamiento de agent chains; borrarlo impediría reproducir la falla en el futuro.
\end{warningbox}

\subsection{Resumen de la sección}
\begin{itemize}
    \item Mediante un flujo con plantilla, un success individual se puede definir como la conclusión de las tres etapas «Plan + Act + Verify».
    \item Registrar la matriz de fallas ofrece una base para el entrenamiento posterior y el ajuste de prompts.
\end{itemize}
\section{Respuesta a incidentes y revisión post mortem}
\subsection{Proceso de atribución de fallas}
Cada vez que un Agente falla, el equipo inicia de inmediato una revisión post mortem: primero extrae la `failure matrix`, que incluye el comando, el archivo que la desencadenó, el entorno de ejecución y el código de retorno; después compara la estrategia actual con las estrategias históricas para determinar si la regresión es del prompt, de la cadena de herramientas o del entorno de ejecución.

\begin{table}[H]
\centering
\begin{tabular}{p{0.3\textwidth} p{0.3\textwidth} p{0.3\textwidth}}
\toprule
Categoría & Condición desencadenante típica & Acción de revisión post mortem \\
\midrule
Falla del prompt & La generación no coincide con el contexto & Rediseñar el prompt como plantilla y agregar nuevos examples \\
\midrule
Error de la cadena de herramientas & El comando o la prueba se cuelga & Revisar la imagen de docker, las versiones de las dependencias y los permisos \\
\midrule
Diferencia en el entorno de evaluación & Pasa en local pero falla en CI & Comparar las variables de entorno y la versión de los datos, y hacer pruebas de dev contra prod \\
\bottomrule
\end{tabular}
\caption{Categorías de fallas y acciones de revisión post mortem.}
\end{table}

\subsection{Difusión de las minutas de revisión post mortem}
El resultado de la revisión post mortem se escribe en la wiki del equipo en formato Markdown, con la fecha, la persona responsable, la conclusión de la revisión y las acciones posteriores, y se envía un recordatorio por el canal de Slack \texttt{\#agent-postmortem}. Cada revisión post mortem también incluye capturas de pantalla clave (como el log de las pruebas o el diff) para ubicar el problema rápidamente en el futuro.

\begin{warningbox}{No dejes que la revisión post mortem se pierda}
Las minutas de la revisión post mortem deben publicarse dentro de las primeras 24 horas; si se tarda demasiado, el escenario se olvida.
\end{warningbox}

\begin{knowledgebox}{Valor de la revisión post mortem}
La revisión post mortem no solo sirve para encontrar errores: convierte las muestras de falla del Agente en material didáctico para entrenar prompts, y transforma el «por qué falló» en un patrón reproducible.
\end{knowledgebox}

\subsection{Resumen de la sección}
\begin{itemize}
    \item La atribución rápida y la revisión post mortem evitan que problemas similares se repitan.
    \item Compartir los resultados de la revisión post mortem con el equipo mejora la madurez general de la ingeniería de Agentes.
\end{itemize}

\section{Verificación de implementación y mejora continua}
\subsection{Lista de verificación habitual}
Graham recomienda ejecutar una lista de verificación antes y después de cada iteración del Agent, para asegurar que el prompt, la cadena de herramientas, los registros y el mecanismo de aprobación estén en condiciones de uso. Los puntos de la lista de verificación incluyen: si se actualizaron los open-source requirements, si se verificaron los permisos de la imagen, si se actualizó la línea base de pruebas, si se sincronizó la notificación de Slack.

\begin{table}[H]
\centering
\begin{tabular}{p{0.32\textwidth} p{0.6\textwidth}}
\toprule
Punto de verificación & Descripción \\
\midrule
Estado del prompt & Asegurar que la plantilla de prompt incluya las muestras de issue más recientes y los guardrails \\
\midrule
Actualización de la cadena de herramientas & Verificar si la imagen de docker, las versiones de dependencias y los permisos de la API están vencidos \\
\midrule
Registros y auditoría & Verificar que el flujo de eventos sea normal y que el audit log pueda reproducirse \\
\midrule
Mecanismo de aprobación & Revisar si el proceso de aprobación de comandos de alto riesgo está completo, incluyendo el emergency override \\
\bottomrule
\end{tabular}
\caption{Puntos de verificación habituales antes y después de la implementación.}
\end{table}

\subsection{Ciclo de mejora continua}
Cada iteración del Agent debe incluir una etapa de revisión: revisar los casos exitosos y las experiencias de fracaso, y escribir los nuevos insight en el FAQ o la documentación, para usarlos en la siguiente ronda de entrenamiento del modelo. Graham enfatiza la «normalización de la mejora del 1%»: cada día mejorar un poco alguno de los elementos entre las instrucciones, los prompts y la cadena de herramientas.

\begin{importantbox}{Pasos pequeños, avance rápido}
Descomponer las actualizaciones del Agent en cambios de 5 minutos, y fusionarlos solo después de que pasen la verificación, evitando en la medida de lo posible los cambios masivos de una sola vez.
\end{importantbox}

\begin{knowledgebox}{Indicadores de mejora}
Usar la «proporción de casos exitosos», la «tasa de finalización de las retrospectivas» y el «número de actualizaciones de documentación» como indicadores de mejora, para cerrar el ciclo entre el trabajo de ingeniería y el de investigación.
\end{knowledgebox}

\subsection{Resumen de la sección}
\begin{itemize}
    \item La lista de verificación garantiza que la arquitectura del Agent esté en buen estado antes y después de cada iteración.
    \item Los pasos pequeños y rápidos junto con los indicadores de mejora permiten que el equipo progrese de manera continua bajo un riesgo controlado.
\end{itemize}

\section{Síntesis y ampliación}
Esta entrevista con Graham puede resumirse en tres temas que avanzan en paralelo: arquitectura, evaluación y gobernanza. Cada uno de estos caminos requiere una práctica de ingeniería intensiva, una evaluación integral y una reflexión sobre seguridad para que los agentes sean ampliamente adoptados en el desarrollo de software de alto riesgo.

\begin{table}[H]
\centering
\begin{tabular}{p{0.34\textwidth} p{0.58\textwidth}}
\toprule
Tema & Puntos clave \\
\midrule
Contexto y motivación & El software está devorando todas las industrias; delegar la capacidad de «escribir código» a los agentes permite que más equipos aborden los problemas con una mentalidad de ingeniería. \\
\midrule
Arquitectura y evaluación & OpenHands ofrece una plataforma de agentes transparente, y SWE-bench Lite/Full ofrece métricas de evaluación de ingeniería con dificultad progresiva. \\
\midrule
Futuro y gobernanza & Se necesitan datos con estilo de agente, más dimensiones de evaluación y mecanismos de seguridad integrados para que los agentes puedan incorporarse con seguridad al flujo de desarrollo convencional. \\
\bottomrule
\end{tabular}
\caption{Los tres elementos centrales de esta charla.}
\end{table}

Lecturas adicionales:
\begin{itemize}
    \item Repositorio de código abierto de OpenHands: \href{https://github.com/OpenHands/OpenHands}{https://github.com/OpenHands/OpenHands}.
    \item Artículo de SWE-bench: Jimenez et al., ``SWE-bench: Can Language Models Resolve Real-World GitHub Issues?'' ICLR 2024.
    \item Blog oficial de All Hands AI y su canal de retroalimentación de la comunidad, enfocados en la seguridad de los agentes y la colaboración humano-máquina.
\end{itemize}

\subsection{Resumen de la sección}
\begin{itemize}
    \item Las tres líneas que enfatiza esta charla —arquitectura, evaluación y gobernanza— deben avanzar en paralelo para llevar a los agentes de forma segura a producción.
    \item Seguir la evolución de OpenHands y SWE-bench ayuda a llevar el desarrollo de software agéntico a la práctica.
\end{itemize}

\end{document}
